% =====================================================================
%  06 模型的评价及推广(建模手)
%  结构:6.1 模型检验 → 6.2 灵敏度分析 → 6.3 优缺点
%
%  占位符约定:\textbf{(待回填:……)} 是给编程手留的数值空位,
%  求解脚本跑完后逐个替换,不要删除。
% =====================================================================
\section{模型的评价及推广}

\subsection{模型检验}

\subsubsection{模型精度检验}

% 待写 + 待新增:精度检验结果表。

\subsubsection{结果合理性检验}

% 待写。

\subsection{灵敏度分析}

% 待写:与模型求解部分(04 对应小节)的敏感性结果保持一致。

\subsection{模型的优缺点}

\textbf{优点:}
% 待写(用 itemize 环境)。
\textbf{缺点:}
% 待写(用 itemize 环境)。
